% from ega1-10-fr.tex
\oldpage[I]{189}

canonique
$f_{\mu\lambda}:(\mathfrak{X},
\mathscr{O}_{\mathfrak{X}}/\mathscr{J}_\lambda)\to
(\mathfrak{X},\mathscr{O}_{\mathfrak{X}}/\mathscr{J}_\mu)$ de préschémas
(usuels) tel que l'on ait $f_\lambda=f_\mu\circ f_{\mu\lambda}$. Les
préschémas
$X_\lambda=(\mathfrak{X},
\mathscr{O}_{\mathfrak{X}}/\mathscr{J}_\lambda)$ et les morphismes
$f_{\mu\lambda}$ constituent donc (en vertu de
\hyperref[I.10.4.8-fr]{10.4.8}) un \emph{système inductif} dans la
catégorie des préschémas formels.
\end{env}

\begin{proposition}[10.6.2]
\label{I.10.6.2-fr}
Avec les notations de (\hyperref[I.10.6.1-fr]{10.6.1}), le préschéma
formel $\mathfrak{X}$ et les morphismes $f_\lambda$ constituent une
limite inductive (T, I, 1.8) du système $(X_\lambda,f_{\mu\lambda})$
dans la catégorie des préschémas formels.
\end{proposition}

\begin{proof}
Soit $\mathfrak{Y}$ un préschéma formel, et pour chaque indice $\lambda$,
soit
\[
  g_\lambda=(\psi_\lambda,\theta_\lambda):X_\lambda\to\mathfrak{Y}
\]
un morphisme, tel que l'on ait
$g_\lambda=g_\mu\circ f_{\mu\lambda}$ pour
$\mathscr{J}_\mu\subset\mathscr{J}_\lambda$. Cette dernière condition
et la définition des $X_\lambda$ entraînent d'abord que tous les
$\psi_\lambda$ sont identiques à une même application continue
$\psi:\mathfrak{X}\to\mathfrak{Y}$ des espaces sous-jacents~; en outre,
les homomorphismes
$\theta_\lambda^\sharp:\psi^*(\mathscr{O}_{\mathfrak{Y}})\to
\mathscr{O}_{X_\lambda}=
\mathscr{O}_{\mathfrak{X}}/\mathscr{J}_\lambda$ forment un
\emph{système projectif} d'homomorphismes de faisceaux d'anneaux. Par
passage à la limite projective, on en déduit donc un homomorphisme
\[
  \omega:\psi^*(\mathscr{O}_{\mathfrak{Y}})\to
  \varprojlim\mathscr{O}_{\mathfrak{X}}/\mathscr{J}_\lambda
  =\mathscr{O}_{\mathfrak{X}},
\]
et il est clair que le morphisme $g=(\psi,\omega^b)$ d'\emph{espaces
annelés} est le \emph{seul} rendant commutatif les diagrammes
\[
  \label{I.10.6.2.1-fr}
  \xymatrix{
    X_\lambda\ar[rr]^{g_\lambda}\ar[rd]_{f_\lambda} &&
    \mathfrak{Y}\\
    & \mathfrak{X}\ar[ru]_g
  }
  \tag{10.6.2.1}
\]
Il reste donc à prouver que $g$ est un morphisme de \emph{préschémas
formels}~; la question étant locale sur $\mathfrak{X}$ et
$\mathfrak{Y}$, on peut donc supposer
$\mathfrak{X}=\operatorname{Spf}(A)$,
$\mathfrak{Y}=\operatorname{Spf}(B)$, $A$ et $B$ étant des anneaux
admissibles, avec
$\mathscr{J}_\lambda=\mathfrak{J}_\lambda^{\Delta}$, où
$(\mathfrak{J}_\lambda)$ est un système fondamental d'idéaux de
définition de $A$ (\hyperref[I.10.3.5-fr]{10.3.5})~; comme
$A=\varprojlim A/\mathfrak{J}_\lambda$, l'existence d'un morphisme de
schémas formels affines $g$ rendant commutatifs les diagrammes
(\hyperref[I.10.6.2.1-fr]{10.6.2.1}) résulte alors de la correspondance
biunivoque (\hyperref[I.10.2.2-fr]{10.2.2}) entre morphismes de schémas
formels affines et homomorphismes continus d'anneaux, et de la définition
de la limite projective. Mais l'unicité de $g$ en tant que morphisme
d'espaces annelés montre qu'il coïncide avec le morphisme noté de même au
début de la démonstration.
\end{proof}

La proposition suivante établit, sous certaines conditions
supplémentaires, l'existence de la limite inductive d'un système inductif
donné de préschémas (usuels) dans la catégorie des préschémas formels~:

\begin{proposition}[10.6.3]
\label{I.10.6.3-fr}
Soient $\mathfrak{X}$ un espace topologique,
$(\mathscr{O}_i,u_{ji})$ un système projectif de faisceaux d'anneaux sur
$\mathfrak{X}$, ayant $\mathbf{N}$ pour ensemble d'indices. Soit
$\mathscr{J}_i$ le noyau de
$u_{0i}:\mathscr{O}_i\to\mathscr{O}_0$. On suppose que~:
\begin{enumerate}
  \item[a)] L'espace annelé $(\mathfrak{X},\mathscr{O}_i)$ est un
    préschéma $X_i$.
  \item[b)] Pour tout $x\in X$ et tout $i$, il existe un voisinage
    ouvert $U_i$ de $x$ dans $\mathfrak{X}$ tel que la restriction
    $\mathscr{J}_i|U_i$ soit nilpotente.
  \item[c)] Les homomorphismes $u_{ji}$ sont surjectifs.
\end{enumerate}

